\section{Diseño del controlador y estimadores}
\subsection{Realimentación de estados estimados}
Ambas arquitecturas utilizan
\begin{equation}\label{eq:control}
u[k]=-K\widehat{\mathbf{x}}[k]+N_r r[k].
\end{equation}
Con los polos de control $0.8\pm j0.25$ y el modelo identificado se obtiene
\begin{equation}
K=\begin{bmatrix}3.0726&9.3389\end{bmatrix},\qquad N_r=0.558625.
\end{equation}
La matriz $A-BK$ es común a los cuatro casos. Una vez que el error de estimación se
vuelve despreciable, ella determina la dinámica del seguimiento.

\subsection{Estimador de predicción}
La estimación disponible para el control corresponde a la predicción previa:
\begin{equation}
\widehat{\mathbf{x}}[k+1]=A\widehat{\mathbf{x}}[k]+Bu[k]
+L_p\bigl(y[k]-C\widehat{\mathbf{x}}[k]\bigr).
\end{equation}
El error evoluciona según $(A-L_pC)e[k]$. Para E1 se adoptan polos
$0.25\pm j0.05$ y para E2, $0.20\pm j0.10$. La ubicación E1 se seleccionó
entre 35 candidatas exigiendo cero abortos con el ruido medido, tolerancia a
un impulso de dos cuentas ADC y un módulo mayor que E2. El diseño produce
\begin{align}
L_{p,E1}&=\begin{bmatrix}-1.2311&0.37767\end{bmatrix}^{T},\\
L_{p,E2}&=\begin{bmatrix}-1.4923&0.47767\end{bmatrix}^{T}.
\end{align}

\subsection{Estimador de estado actual}
El estimador actual corrige la predicción con la muestra presente antes de
calcular~\eqref{eq:control}:
\begin{align}
\mathbf z[k]&=A\widehat{\mathbf{x}}[k-1]+Bu[k-1],\\
\widehat{\mathbf{x}}[k]&=\mathbf z[k]+L_a\bigl(y[k]-C\mathbf z[k]\bigr).
\end{align}
Su error corregido evoluciona con $(I-L_aC)A$. Para los mismos polos se obtiene
\begin{align}
L_{a,E1}&=\begin{bmatrix}-1.5348&-0.072302\end{bmatrix}^{T},\\
L_{a,E2}&=\begin{bmatrix}-1.8605&0.17515\end{bmatrix}^{T}.
\end{align}
Las ordenadas iniciales de predictor y actual no son directamente equivalentes:
el actual ya incorporó $y[k]$ en la muestra representada, mientras el predictor
aplicará esa innovación a la estimación siguiente.

\subsection{Límites y cuantización}
La solicitud se limita a $[-115,107]$ mV. La cadena empleada produce pasos
nominales de \SI{15.37}{\milli\volt} para la rama positiva y
\SI{10.56}{\milli\volt} para la negativa. Por tanto, diferencias pequeñas en
la acción calculada pueden desaparecer al convertirla en códigos del DAC. Esta
no linealidad se incluye en la simulación alineada con la campaña experimental.
